\section{Why On-Policy RL Can Lose Modes and Verified Replay Retains Them}
\label{app:theory}

\noindent\emph{Main-body link.} This appendix supplies the formal claims
invoked in Sections~\ref{sec:method} and~\ref{sec:results}.

This appendix separates two questions that are easy to conflate: whether a
binary-reward group update preserves ratios among already correct modes, and
whether verified replay preserves modes the policy has produced
and the validator has accepted. The first admits a generic-collapse result in a
finite model with independently trained categorical logits. The second necessarily requires a coverage
assumption: the replay term directly protects only represented bank modes;
it does not by itself guarantee retention of modes never sampled or constructed. We therefore
prove support preservation first and state full no-collapse as a coverage
corollary. Accordingly, the positive theorem explains the stability mechanism
after discovery; it does not prove that a ReplayMaxRL policy will establish its own
coverage premise in a language model.

\paragraph{How to read the proof chain.}
The group-expectation lemmas first reduce the fresh task update to a positive
multiple of the correctness gradient. The collapse theorem then separates
that scalar learning clock from the within-correct-mode dynamics. The shared
correctness theorem compares those dynamics at matched accuracy, and
Corollary~\ref{cor:collapse-error-rate} quantifies their eventual loss.
Lemmas~\ref{lem:entropy-scope} and~\ref{lem:sampled-score-bound} separate
exact entropy regularization from bounded sampled surprisal.
Lemma~\ref{lem:replay-gradient-availability} distinguishes a fresh task signal
from a replay signal. The retention theorem and its fixed-bank limit then
identify exactly which stored modes survive. The natural-gradient comparison
in Section~\ref{app:theory-natural-gradient} makes the role of update geometry
explicit, while Section~\ref{app:theory-entropy-comparison} relates exact entropy
to replay through their distinct relative-entropy objectives.
Lemma~\ref{lem:shared-exemplar-retention} then states the energy condition
needed for complete exemplars in a shared model. The remaining sections extend
that condition to controlled stochastic updates and changing banks, and derive
sharp per-mode certificates and valid finite-sample interpretations
(Sections~\ref{app:theory-optimizer}--\ref{app:theory-survival-certificates}).
Each extension states its additional assumptions and the established
literature tools used in its proof.

Here $G$ counts training rollouts, $K$ counts evaluation draws, $m,n$ count
correct and incorrect categories, and $k$ is the bank size. In the mean-flow results, $t$ is
continuous optimization time and $\tau$ is a reparameterized learning clock.
The later stochastic and admission results use explicitly defined update or
opportunity indices. None is a wall-clock or GPU-runtime prediction. ``Collapse'' refers to
limiting concentration: finite softmax logits already give positive
probabilities at every finite time. A uniform lower bound for all future
times is the stronger retention property proved below.

\subsection{Categorical reduction and the GRPO mean update}
\label{app:theory-grpo-mean}

Fix one prompt. Let the finite response categories be
$\mathcal A=\mathcal C\mathbin{\dot\cup}\mathcal N$, where the $m\ge2$
categories in $\mathcal C$ are correct execution modes and the $n\ge1$
categories in $\mathcal N$ are incorrect. Let
$p_a=\operatorname{softmax}(z)_a$, $P=\sum_{c\in\mathcal C}p_c$, and
$q_c=p_c/P$. Thus $q$ is exactly the conditional correct-mode distribution
$p_\theta^+(\cdot\mid x)$ of Section~\ref{sec:collapse}, and $P$ is the correct
mass $P_\theta^+(x)$.

Draw a group $A_1,\ldots,A_G\stackrel{\mathrm{iid}}{\sim}p$ with $G\ge2$,
write $R_i=\1[A_i\in\mathcal C]$, and let $M=\sum_iR_i$. Both GRPO variants
can be written, at the on-policy point where the PPO ratio equals one, as
\begin{equation}
  \widehat g_G
  =\frac1G\sum_{i=1}^G
   w(M)\left(R_i-\frac{M}{G}\right)\nabla_z\log p_{A_i}.
  \label{eq:group-estimator}
\end{equation}
For Dr.GRPO, $w(M)=1$ after absorbing the shared positive length scale. For
binary-reward GRPO, $w(M)$ is the reciprocal group reward standard deviation
on mixed groups. We absorb the same fixed response-length normalizer in both
variants; response-dependent row normalization is outside this reduction.
In either case $0<w(\ell)<\infty$ for $\ell\in\{1,\ldots,G-1\}$.
We define the entire update as zero on all-equal groups, without evaluating an
undefined reciprocal standard deviation; likewise, $w(M)M(G-M)$ below is
defined as zero when $M=0$ or $M=G$. We take explicit reference KL to be zero, as
in the experiments.

\begin{lemma}[A group baseline does not remove frequency amplification]
\label{lem:grpo-mean}
For $0<P<1$, the expected group update is
\begin{equation}
  \E[\widehat g_G]
  =c_G(P)\nabla_z P,
  \qquad
  c_G(P)=
  \frac{\E\!\big[w(M)M(G-M)\big]}{G^2P(1-P)}>0.
  \label{eq:expected-group-gradient}
\end{equation}
For Dr.GRPO, $c_G(P)=(G-1)/G$. Hence reward centering and binary reward
standardization change the speed, but not the expected direction, of the
binary-reward policy gradient.
\end{lemma}

\begin{proof}
Condition on the complete binary reward vector, whose sum is $M$. The score
identities give
\[
  \E[\nabla_z\log p_{A_i}\mid R_i=1]=\nabla_z\log P,
  \qquad
  \E[\nabla_z\log p_{A_i}\mid R_i=0]
  =\nabla_z\log(1-P).
\]
There are $M$ correct rows with centered reward $(G-M)/G$ and $G-M$
incorrect rows with centered reward $-M/G$. Therefore the conditional
expectation of the unaveraged sum in~\eqref{eq:group-estimator} is
\[
  w(M)\frac{M(G-M)}{G}
  \left(\nabla_z\log P-\nabla_z\log(1-P)\right)
  =w(M)\frac{M(G-M)}{GP(1-P)}\nabla_zP.
\]
The outer factor $1/G$ and expectation over
$M\sim\operatorname{Binomial}(G,P)$ yield~\eqref{eq:expected-group-gradient}.
For $w\equiv1$,
$\E[M(G-M)]=G(G-1)P(1-P)$, proving the Dr.GRPO value.
\end{proof}

The binomial identity also gives the denominator-free expression
\[
 c_G(P)=\frac{G-1}{G}\sum_{j=0}^{G-2}\binom{G-2}{j}
 w(j+1)P^j(1-P)^{G-2-j}.
\]
The binomial weights sum to one, so $c_G$ lies between
$(G-1)\min_j w(j+1)/G$ and $(G-1)\max_j w(j+1)/G$.
This establishes the smooth, positive extension to $P=0,1$ used in the
subsequent flow and convergence arguments.

The next lemma specializes the practical-estimator analysis of
\citet[Section~4.3 and Appendix~D, Theorem~5]{tajwar2026maxrl}
to our categorical notation; we include the derivation for completeness.

\begin{lemma}[Binary MaxRL has the same correctness-gradient direction]
\label{lem:maxrl-mean}
Let MaxRL assign
$A_i^{\mathrm{MaxRL}}=GR_i/M-1$ when $M>0$ and zero when $M=0$, as in
Section~\ref{sec:maxrl-method}. At the on-policy point, its expected update is
\begin{equation}
  \E[\widehat g_G^{\mathrm{MaxRL}}]
  =c_G^{\mathrm{MaxRL}}(P)\nabla_zP,
  \qquad
  c_G^{\mathrm{MaxRL}}(P)
  =\frac{1-(1-P)^{G-1}}{P}>0,
  \label{eq:maxrl-expected-gradient}
\end{equation}
with continuous endpoint values
$c_G^{\mathrm{MaxRL}}(0)=G-1$ and
$c_G^{\mathrm{MaxRL}}(1)=1$.
\end{lemma}

\begin{proof}
Condition on $M=m>0$. Each successful row has advantage $(G-m)/m$ and each
failed row has advantage $-1$. Using the same conditional score identities as
above, the expected unaveraged sum is
\[
  (G-m)\!\left(\nabla_z\log P-\nabla_z\log(1-P)\right)
  =\frac{G-m}{P(1-P)}\nabla_zP.
\]
After the outer factor $1/G$ and averaging over
$M\sim\operatorname{Binomial}(G,P)$,
\[
 \E[(G-M)\1\{M>0\}]
 =G(1-P)-G\Pr(M=0)
 =G\big[(1-P)-(1-P)^G\big].
\]
Substitution gives Equation~\eqref{eq:maxrl-expected-gradient}. Its endpoint
values follow by continuity (or the finite geometric-series identity), which
also gives the stated continuous extension at both boundaries.
\end{proof}

The same finite-series identity gives the potential used below:
\begin{equation}
 \Psi_{\mathrm{MaxRL}}(P)
 :=\int_0^P c_G^{\mathrm{MaxRL}}(s)\,ds
 =\sum_{k=1}^{G-1}\frac{1-(1-P)^k}{k}.
 \label{eq:maxrl-potential-finite-sum}
\end{equation}
In particular, $\Psi_{\mathrm{MaxRL}}(1)=\sum_{k=1}^{G-1}1/k$.
For binary rewards, best@\(k\) equals pass@\(k\), so this is a harmonic
mixture of sampling objectives for $k=1,\ldots,G-1$. The upper limit is
$G-1$, not $G$: the implemented centered estimator drops the entire update
on all-failure groups. Subtracting an unconditional score baseline even on
those groups instead gives the order-$G$ estimator. This distinction concerns
the exact finite-group expectation at the on-policy point; it does not assert
unbiasedness of subsequent clipped or adaptive-optimizer updates.

\subsection{Winner-take-all collapse under GRPO, Dr.GRPO, and MaxRL}
\label{app:theory-collapse}

Prior work studies objective-induced outcome collapse and inverse-probability
correction \citep{sinha2026expected}; UCPO analyzes correctness-indifferent collapse
through sampling bias and gradient amplification \citep[Section~4]{lochab2026ucpo}.
Here we connect our group estimators to a standard categorical flow and then study
persistent replay over execution-defined keys.
Consider the infinitesimal expected-update flow
$\dot z=\E[\widehat g_G]=c_G(P)\nabla_zP$. This is the clean mean-field limit
of fresh on-policy groups and the first-order PPO update. By Lemmas~\ref{lem:grpo-mean} and~\ref{lem:maxrl-mean}, it covers GRPO, Dr.GRPO, and binary MaxRL with their respective positive functions $c_G$. It deliberately
removes sampling noise; the result below therefore identifies a structural
instability rather than blaming finite batches.

The conditional correct-mode flow derived below is a standard
frequency-dependent replicator system with fitness $f_c(q)=q_c$
\citep{hofbauer1998evolutionary}; the following
theorem is not claimed as a new replicator-dynamics result. Its role is to show
that the expected GRPO, Dr.GRPO, and binary-MaxRL updates reduce to that
familiar winner-take-all flow under the paper's categorical assumptions.

\begin{theorem}[Standard replicator specialization for binary-reward RL]
\label{thm:grpo-collapse}
Start from finite logits with $0<P(0)<1$. If the conditional correct-mode
distribution $q(0)$ has a unique largest coordinate $c^\star$, then
\[
  P(t)\longrightarrow1,
  \qquad
  q_{c^\star}(t)\longrightarrow1,
  \qquad
  q_c(t)\longrightarrow0\quad(c\ne c^\star).
\]
Thus these GRPO, Dr.GRPO, and binary MaxRL mean flows converge to a single
correct execution mode for generic initial conditions, even though every mode
in $\mathcal C$ has exactly the same reward.
\end{theorem}

\begin{proof}
For a correct mode $c$ and incorrect category $a$,
\[
  \dot z_c=c_G(P)p_c(1-P),
  \qquad
  \dot z_a=-c_G(P)p_aP.
\]
Also $\dot P=c_G(P)\lVert\nabla_zP\rVert_2^2>0$. If $P$ converged to an
interior value, continuity and positivity of $c_G$, together with
\[
 \lVert\nabla_zP\rVert_2^2
 \ge P^2(1-P)^2\left(\frac1m+\frac1n\right),
\]
would keep $\dot P$ bounded away from zero, a contradiction. Hence $P\to1$.

Define the effective optimization time
\[
  \tau(t)=\int_0^t c_G(P(s))P(s)(1-P(s))\,ds.
\]
If $\tau(\infty)<\infty$, every correct and incorrect logit would have finite
total variation, because the absolute velocity of each is at most $\dot\tau$.
All logits would then converge to finite values, contradicting $P\to1$.
Therefore $\tau(t)\to\infty$, leaving infinite effective time for the
correct-simplex dynamics below.

On the correct simplex, changing from $t$ to $\tau$ gives the replicator flow
\begin{equation}
  \frac{dq_c}{d\tau}
  =q_c\left(q_c-\sum_{d\in\mathcal C}q_d^2\right),
  \qquad
  \frac{d}{d\tau}\log\frac{q_c}{q_d}=q_c-q_d.
  \label{eq:grpo-replicator}
\end{equation}
The initially unique maximum remains the unique maximum. For $c\ne c^\star$,
set $\xi_c=q_c/q_{c^\star}$, so $\xi_c(0)<1$ and $\xi_c$ is nonincreasing.
Since $q_{c^\star}\ge1/m$,
\[
 \frac{d\log\xi_c}{d\tau}
 =-q_{c^\star}(1-\xi_c)
 \le-\frac{1-\xi_c(0)}{m}.
\]
Hence $\xi_c(\tau)\le\xi_c(0)\exp[-(1-\xi_c(0))\tau/m]\to0$.
Every competing mode vanishes relative to $c^\star$; normalization then
yields $q_{c^\star}\to1$, establishing the claimed single-mode limit.
\end{proof}

\begin{corollary}[Ties are nongeneric]
\label{cor:grpo-ties}
If exactly $k<m$ correct modes tie for the largest initial probability, symmetry
preserves the tie and the flow converges to the uniform distribution on those
$k$ modes while eliminating the rest. Only the exactly uniform initialization preserves all $m$ correct modes
in the limit; a tie among $1<k<m$ maxima yields partial collapse onto those
$k$ modes. Any perturbation creating a unique maximum gives single-mode
collapse by Theorem~\ref{thm:grpo-collapse}.
\end{corollary}

Finite groups also introduce sampling effects: a mode absent from
a group has no sampled response of its own carrying a score term,
all-correct groups have zero task advantage, and sampling noise can break exact
ties. Softmax normalization and shared parameters can still change an
unsampled mode's probability. These facts are not needed for
the theorem.

\subsection{How model parameterization changes collapse}
\label{app:theory-size}

The cross-scale losses raise a further question: can a larger model learn
correctness while preserving more correct alternatives? The preceding theorem
already allows every logit to vary independently, so representational capacity
alone cannot answer this question. We give a conditional mechanism: directions
that raise correct logits together can improve correctness with less
amplification of differences among correct modes.

For a smooth parameterization $z=z(\theta)$, write
$J_\theta=\partial z/\partial\theta$ and
$\mathsf K_\theta=J_\theta J_\theta^\top$. The chain rule extends
Lemmas~\ref{lem:grpo-mean} and~\ref{lem:maxrl-mean} to Euclidean parameter
gradient flow:
\begin{equation}
 \dot z=c_G(P)\mathsf K_\theta\nabla_zP,
 \qquad
 \frac{d}{dt}\log\frac{q_c}{q_d}
 =c_G(P)(\mathbf e_c-\mathbf e_d)^\top\mathsf K_\theta\nabla_zP,
 \label{eq:parameterized-collapse}
\end{equation}
where $\mathbf e_c$ is a coordinate basis vector. Thus the logit-gradient geometry,
not the number of columns of $J_\theta$ alone, governs mode amplification.
The independent-logit theorem has $\mathsf K_\theta=I$.

To isolate shared correctness learning, let $v=\1_{\mathcal C}$ be the
vector that is one on correct categories and zero on incorrect categories,
and consider the constant geometry
\begin{equation}
 \mathsf K_\lambda=I+\lambda vv^\top,\qquad \lambda\ge0.
 \label{eq:shared-correctness-kernel}
\end{equation}
The added direction moves all correct logits equally, so it improves their
mass without directly changing their ratios. For example,
$z=u+\alpha\sum_{j=1}^s h_jv$, with trainable $u$ and scalar $h_j$, realizes
this geometry with $\lambda=s\alpha^2$. More such coordinates strengthen the
shared update if $\alpha$ is fixed; normalizing $\alpha=1/\sqrt{s}$ removes
that dependence. These coordinates are redundant for expressivity and assume
perfect correctness alignment. They describe an optimization mechanism, not
a feature that an arbitrary larger language model is guaranteed to possess.

\begin{theorem}[Shared correctness learning preserves more sampled modes at matched accuracy]
\label{thm:shared-correctness}
Compare the flows~\eqref{eq:parameterized-collapse} with
$\mathsf K_\theta=\mathsf K_\lambda$, the same finite initial logits, and any
of the three functions $c_G$ above. Fix a target correct mass
$P_\star\in(P(0),1)$. At the first time each flow reaches $P_\star$,
$\texttt{distinct@}K$ is nondecreasing in $\lambda$ for every integer $K\ge1$,
while $\texttt{pass@}K$ is identical. The increase is strict for $K\ge2$ when
$q(0)$ is nonuniform. In addition, every correct mode satisfies
\begin{equation}
 q_c^{(\lambda)}(P_\star)
 \ge q_c(0)\exp\!\left(-\frac{L}{\lambda+1/m+1/n}\right),
 \quad
 L=\log\frac{P_\star}{1-P_\star}-\log\frac{P(0)}{1-P(0)}.
 \label{eq:shared-correctness-bound}
\end{equation}
Here $m,n$ count output categories, not model parameters. For every fixed
finite $\lambda$, a unique initial maximum still implies eventual
single-mode collapse as $t\to\infty$.
\end{theorem}

\begin{proof}
Let $r_a=p_a/(1-P)$ for $a\in\mathcal N$, and write
$S_q=\sum_cq_c^2$, $S_r=\sum_ar_a^2$. In the effective time
$d\tau=c_G(P)P(1-P)\,dt$, the logit velocities are
$dz_c/d\tau=q_c+\lambda$ and $dz_a/d\tau=-r_a$. Hence
\begin{equation}
 \frac{dq_c}{d\tau}=q_c(q_c-S_q),\qquad
 \frac{dr_a}{d\tau}=-r_a(r_a-S_r),\qquad
 \frac{d}{d\tau}\log\frac{P}{1-P}=S_q+S_r+\lambda.
 \label{eq:shared-correctness-clock}
\end{equation}
The paths $q(\tau),r(\tau)$ therefore do not depend on $\lambda$.
Also $\dot P=c_G(P)P^2(1-P)^2(S_q+S_r+\lambda)>0$; the same interior-limit
argument as in Theorem~\ref{thm:grpo-collapse} gives $P\to1$.
At $P_\star$, the effective time $\tau_\lambda$ uniquely solves
\[
 L=\lambda\tau_\lambda+
   \int_0^{\tau_\lambda}(S_q(u)+S_r(u))\,du.
\]
The right-hand side increases strictly with both its positive time argument
and $\lambda$, so $\tau_\lambda$ decreases strictly with $\lambda$.
Since $S_q\ge1/m$ and $S_r\ge1/n$,
$\tau_\lambda\le L/(\lambda+1/m+1/n)$. Integrating
$d\log q_c/d\tau=q_c-S_q\ge-1$ over this bounded interval gives the probability
floor in Equation~\eqref{eq:shared-correctness-bound}.

At fixed $P_\star$, define
$\mathcal D_K(q)=\sum_c[1-(1-P_\star q_c)^K]$, the expected number of correct modes in
$K$ independent draws. For
$a(u)=KP_\star(1-P_\star u)^{K-1}$, the replicator equation gives
\[
 \frac{d\mathcal D_K}{d\tau}
 =\frac12\sum_{c,d}q_cq_d(q_c-q_d)\bigl(a(q_c)-a(q_d)\bigr)\le0.
\]
The inequality holds because $a$ is nonincreasing; it is strict for $K\ge2$
and nonuniform $q$. Thus the smaller $\tau_\lambda$ gives larger
$\texttt{distinct@}K$, whereas
$\texttt{pass@}K=1-(1-P_\star)^K$ is unchanged. Finally,
$S_q+S_r+\lambda\le2+\lambda$ and $P\to1$ imply $\tau\to\infty$.
The unchanged correct-mode replicator flow then gives eventual collapse
under a unique initial maximum.
\end{proof}

\begin{corollary}[Collapse rate as correctness approaches one]
\label{cor:collapse-error-rate}
Under Theorem~\ref{thm:shared-correctness}, fix finite $\lambda\ge0$
and suppose the initial correct-mode maximum is unique. For every other
correct mode $c$,
\begin{equation}
 \lim_{P\uparrow1}\frac{\log q_c^{(\lambda)}(P)}{\log(1-P)}
 =\frac{1}{\lambda+1+1/n}.
 \label{eq:collapse-error-exponent}
\end{equation}
Thus $q_c^{(\lambda)}(P)=(1-P)^{1/(\lambda+1+1/n)+o(1)}$:
larger fixed $\lambda$ reduces the decay exponent, but every minority mode
still vanishes. The exponent describes the limiting ratio of logarithms;
finite-accuracy probabilities also depend on the initial distribution.
\end{corollary}

\begin{proof}
Theorem~\ref{thm:shared-correctness} gives $q\to\mathbf e_{c^\star}$,
so $S_q\to1$. The incorrect conditional flow in
Equation~\eqref{eq:shared-correctness-clock} preserves ordering and makes
$r_{\min}$ nondecreasing. For $h=\log(r_{\max}/r_{\min})$,
\[
 h'=-(r_{\max}-r_{\min})
 =-r_{\min}(e^h-1)\le-r_{\min}(0)h.
\]
Hence $r\to(1/n,\ldots,1/n)$ and $S_r\to1/n$, including trivially $n=1$.
Integrating $d\log q_c/d\tau=q_c-S_q$ and the correctness-clock equation
and taking time averages yields
\[
 \frac{\log q_c(\tau)}{\tau}\to-1,
 \qquad
 \frac{\log(P(\tau)/(1-P(\tau)))}{\tau}
 \to\lambda+1+1/n.
\]
Since $\log(1-P)=\log P-\log(P/(1-P))$ and $\log P\to0$, their ratio proves
Equation~\eqref{eq:collapse-error-exponent}.
\end{proof}

\paragraph{Correctness alone does not imply collapse.}
The independent-logit component matters. With purely shared geometry
$\mathsf K=\lambda vv^\top$, $\lambda>0$, all correct logits move equally,
so $q(t)=q(0)$ while
$\dot P=\lambda c_G(P)P^2(1-P)^2>0$ and $P\to1$.
For $\mathsf K=aI+bvv^\top$ with fixed $a>0$ and $b\ge0$, time rescaling
instead recovers Theorem~\ref{thm:shared-correctness} with $\lambda=b/a$.
Thus persistent mode-specific amplification, not increasing correctness by
itself, drives this collapse result.

\paragraph{Why parameter count alone is insufficient.}
For an integer $h\ge1$, replace each independent logit by
$z_a=h^{-\gamma}\sum_{j=1}^h\theta_{aj}$, initializing
$\theta_{aj}(0)=h^{\gamma-1}z_a(0)$. Every added parameter affects the logits,
and every model represents the same categorical policies, but
$\mathsf K=h^{1-2\gamma}I$. Its policy follows exactly the original collapse
trajectory with time multiplied by $h^{1-2\gamma}$. Increasing the parameter
count therefore makes collapse faster, unchanged, or slower for
$\gamma=0,1/2,1$, respectively. A size claim requires assumptions about how
additional parameters change the update geometry, not just their number.

\paragraph{What this explains and what remains untested.}
Theorem~\ref{thm:shared-correctness} connects a specified kind of shared
learning to the paper's sampling metrics at matched accuracy. It also
clarifies why a slower loss can coexist with eventual collapse, leaving a
role for replay. It does not identify the cause of the empirical cross-scale
pattern: those runs start from different pretrained distributions, use
different learning-rate schedules, and compare equal training passes rather
than matched accuracy. Falcon also changes model family. Neither the assumed
geometry nor its dependence on size is measured here. Establishing that link
would require controlled within-family scaling and measurements of shared
versus mode-specific updates. The result retains prompt isolation, exact
mean updates, and zero reference KL; it does not cover shared-prompt
interference, adaptive optimizer effects, or finite-step stochastic training.

